\section{Conclusion}
\label{sec:conclusion}

% TODO(results): once numbers are final, tighten the evidence sentence below with
% the headline gains (currently qualitative to avoid unsupported claims).

We studied attacker detection in LLM-based multi-agent systems and identified a
temporal blind spot in state-of-the-art graph-based defenses: by mean-pooling an
agent's multi-turn dialogue before reasoning over the utterance graph, detection
depends on average content and ignores behavioral dynamics. We turned this blind
spot into both a threat and a remedy. Our \emph{escalation attack} starts benign
and intensifies over rounds, saturating the aggregated feature so that
aggregation-based detectors fail while the attack still propagates. Our defense,
\textbf{T-GAT}, replaces mean-pooling with a temporal edge encoder that combines
self-attention, temporal convolution, and cross-turn statistics, so that
detection keys on \emph{how} an agent changes rather than only its average
content. Experiments on tool-attack scenarios across multiple LLM backbones and
topologies show that the escalation attack substantially degrades the baseline,
while T-GAT restores detection and lowers attack success rate under both
escalation and continuous attacks, all while preserving the inductive
transferability of graph-based detectors.

\paragraph{Limitations and future work.}
T-GAT is a \emph{reactive} defense: like the detect-and-remediate paradigm it
builds on, it curbs the spread of malicious content but cannot prevent the initial
compromise, and its temporal evidence is bounded by the short dialogue horizon
typical of current MAS. Our evaluation also focuses on tool-attack scenarios and
on attackers whose behavior leaves a temporal trace; a defense-aware adversary
that flattens its trajectory to mimic honest dynamics remains an open challenge.
Promising next steps are proactive prevention rather than post-hoc remediation,
detection over longer interaction horizons, and an explicit study of adaptive,
detection-aware attacks across a broader range of attack families.
